% ============================================================================
% STUNDENLEISTUNG: Mi casa
% ============================================================================
% Sequenz 3, Block d: Stundenleistung (Test)
% 20 Punkte gesamt, 25 Minuten Bearbeitungszeit
%
% Aufgabe 1: Raume zuordnen (4P)
% Aufgabe 2: hay + Mobel (4P)
% Aufgabe 3: Farben + Adjektivanpassung (6P)
% Aufgabe 4: Zimmer beschreiben (6P)
% ============================================================================

\documentclass[11pt, a4paper]{article}

% ============================================================================
% SW-MODUS TOGGLE
% ============================================================================
\newif\ifswmodus
\swmodusfalse

% ============================================================================
% PAKETE
% ============================================================================

\usepackage[utf8]{inputenc}
\usepackage[T1]{fontenc}
\usepackage[ngerman]{babel}
\usepackage{helvet}
\renewcommand{\familydefault}{\sfdefault}

\usepackage[margin=2cm]{geometry}
\usepackage{enumitem}
\usepackage{tabularx}
\usepackage{booktabs}
\usepackage{multirow}
\usepackage{pgffor}

\usepackage{xcolor}
\usepackage{framed}
\usepackage{tcolorbox}
\tcbuselibrary{skins}

\usepackage{fancyhdr}
\usepackage{lastpage}
\usepackage{needspace}

% ============================================================================
% FARBEN
% ============================================================================

\definecolor{spanisch-color}{HTML}{c41e3a}
\definecolor{grau-color}{HTML}{888888}
\definecolor{hellgrau-color}{HTML}{f5f5f5}
\definecolor{orange-color}{HTML}{b35c00}

\definecolor{sw-dark}{HTML}{000000}
\definecolor{sw-gray}{HTML}{666666}
\definecolor{sw-white}{HTML}{ffffff}

\ifswmodus
    \colorlet{spanisch}{sw-dark}
    \colorlet{grau}{sw-gray}
    \colorlet{hellgrau}{sw-white}
    \colorlet{orange}{sw-dark}
\else
    \colorlet{spanisch}{spanisch-color}
    \colorlet{grau}{grau-color}
    \colorlet{hellgrau}{hellgrau-color}
    \colorlet{orange}{orange-color}
\fi

\newcommand{\fachfarbe}{spanisch}

% ============================================================================
% VARIABLEN
% ============================================================================

\newcommand{\thema}{Mi casa -- Stundenleistung}
\newcommand{\sequenz}{03}
\newcommand{\block}{d}
\newcommand{\fach}{Spanisch}
\newcommand{\klasse}{7}
\newcommand{\maxpunkte}{20}
\newcommand{\dauer}{25 Minuten}
\newcommand{\datum}{\today}

% ============================================================================
% KOPF- UND FUSSZEILE
% ============================================================================

\pagestyle{fancy}
\fancyhf{}
\renewcommand{\headrulewidth}{0.4pt}
\renewcommand{\footrulewidth}{0.4pt}

\lhead{\fach{} -- Klasse \klasse}
\chead{\textbf{SL-\sequenz\block{} (Stundenleistung)}}
\rhead{\dauer}

\lfoot{}
\cfoot{}
\rfoot{Seite \thepage\ von \pageref{LastPage}}

% ============================================================================
% BEFEHLE
% ============================================================================

\newcommand{\punktebox}[1]{\hfill\fbox{\small \_/#1}}
\newcommand{\luecke}[1]{\rule{#1}{0.4pt}}

\newtcolorbox{vokabelkasten}{
    colback=white,
    colframe=\fachfarbe,
    boxrule=0.5pt,
    arc=0pt,
    left=5pt, right=5pt, top=5pt, bottom=5pt
}

\newtcolorbox{hinweis}{
    colback=white,
    colframe=orange,
    colbacktitle=white,
    coltitle=orange,
    boxrule=1pt,
    arc=0pt,
    left=5pt, right=5pt, top=5pt, bottom=5pt,
    fonttitle=\bfseries,
    attach boxed title to top left={yshift=-2mm, xshift=5mm},
    boxed title style={boxrule=0pt, colback=white},
    title=Hinweis
}

\newenvironment{aufgabe}[1]{%
    \needspace{5\baselineskip}%
    \nopagebreak[4]%
    \vspace{10pt}
    \noindent\textbf{\textcolor{\fachfarbe}{Aufgabe #1}}
    \nopagebreak[4]%
    \vspace{5pt}
    \nopagebreak[4]%
    \begin{list}{}{\leftmargin=0pt}
    \item[]
}{%
    \end{list}
}

\newcommand{\antwortzeilen}[1]{%
    \vspace{2pt}
    \foreach \n in {1,...,#1}{%
        \noindent\rule{\textwidth}{0.4pt}\vspace{8pt}
    }
}

\newcommand{\es}[1]{\textcolor{\fachfarbe}{\textbf{#1}}}

% ============================================================================
% DOKUMENT
% ============================================================================

\begin{document}

% ============================================================================
% KOPFBEREICH
% ============================================================================

\begin{center}
    {\Large\bfseries\textcolor{\fachfarbe}{Stundenleistung: \thema}}
\end{center}

\vspace{5pt}

\noindent
\begin{tabularx}{\textwidth}{|X|X|X|X|}
\hline
\textbf{Name:} \luecke{3cm} &
\textbf{Klasse:} \klasse &
\textbf{Datum:} \luecke{2cm} &
\textbf{Zeit:} \dauer \\
\hline
\end{tabularx}

\vspace{10pt}

\begin{hinweis}
Lies alle Aufgaben zuerst durch. Schreibe deutlich! Du hast \dauer{} Zeit.
\end{hinweis}

\vspace{15pt}

% ============================================================================
% AUFGABE 1: Raume zuordnen (4 Punkte)
% ============================================================================

\begin{aufgabe}{1: Raume zuordnen}
Ordne die spanischen Raume den deutschen Ubersetzungen zu. Schreibe den passenden Buchstaben in die Lucke. \punktebox{4}
\end{aufgabe}

\begin{center}
\begin{tabular}{lcp{6cm}}
1. \es{la cocina} & = \luecke{1.5cm} & a) das Schlafzimmer \\[0.8em]
2. \es{el dormitorio} & = \luecke{1.5cm} & b) das Badezimmer \\[0.8em]
3. \es{el salon} & = \luecke{1.5cm} & c) die Kuche \\[0.8em]
4. \es{el bano} & = \luecke{1.5cm} & d) das Wohnzimmer \\
\end{tabular}
\end{center}

\vspace{15pt}

% ============================================================================
% AUFGABE 2: hay + Mobel (4 Punkte)
% ============================================================================

\begin{aufgabe}{2: hay + Mobel}
Erganze die Satze mit \es{hay} und einem passenden Mobel aus dem Kasten. \punktebox{4}
\end{aufgabe}

\begin{vokabelkasten}
\textit{Mobel:} una cama | un sofa | una mesa | un armario
\end{vokabelkasten}

\vspace{0.5em}

\noindent
a) En la cocina \luecke{4cm} \luecke{3cm}. \\[1em]
b) En el dormitorio \luecke{4cm} \luecke{3cm}. \\[1em]
c) En el salon \luecke{4cm} \luecke{3cm}. \\[1em]
d) En el dormitorio tambien \luecke{4cm} \luecke{3cm}. \\

\vspace{15pt}

% ============================================================================
% AUFGABE 3: Farben + Adjektivanpassung (6 Punkte)
% ============================================================================

\begin{aufgabe}{3: Farben + Adjektivanpassung}
Erganze die Farbadjektive. Achte auf die richtige Endung! \punktebox{6}
\end{aufgabe}

\begin{vokabelkasten}
\textit{Farben:} verde | rojo/roja | blanco/blanca | marron/marrones | azul | grande
\end{vokabelkasten}

\vspace{0.5em}

\noindent
a) El sofa es \luecke{3cm} (grun). \\[1em]
b) La silla es \luecke{3cm} (rot). \\[1em]
c) Las mesas son \luecke{3cm} (weiß). \\[1em]
d) Los armarios son \luecke{3cm} (braun). \\[1em]
e) La cama es \luecke{3cm} (blau). \\[1em]
f) El jardin es \luecke{3cm} (groß). \\

\vspace{15pt}

% ============================================================================
% AUFGABE 4: Zimmer beschreiben (6 Punkte)
% ============================================================================

\begin{aufgabe}{4: Zimmer beschreiben}
Beschreibe ein Zimmer deiner Wahl auf Spanisch. Schreibe mindestens 4 Satze. Verwende \es{hay}, Mobel und Farben. \punktebox{6}
\end{aufgabe}

\begin{hinweis}
Beispiel: En mi dormitorio hay una cama blanca. La cama es grande...
\end{hinweis}

\vspace{0.5em}

\noindent\textbf{Mi habitacion favorita:} \luecke{4cm} (Raumname auf Spanisch)

\antwortzeilen{7}

% ============================================================================
% PUNKTEUBERSICHT
% ============================================================================

\vspace{20pt}
\hrule
\vspace{10pt}

\noindent
\begin{tabularx}{\textwidth}{|X|c|c|c|c||c|}
\hline
& \textbf{Aufg. 1} & \textbf{Aufg. 2} & \textbf{Aufg. 3} & \textbf{Aufg. 4} & \textbf{Gesamt} \\
\hline
\textbf{Punkte} & \_/4 & \_/4 & \_/6 & \_/6 & \_/\maxpunkte \\
\hline
\end{tabularx}

\vspace{10pt}

\begin{center}
\textbf{Note:} \luecke{2cm}
\end{center}

\end{document}
